\chapter{Justificaci├│n}
\label{chap:justificacion}

\noindent La prevalencia del comportamiento sedentario (CS) en la sociedad contemporánea es un fenómeno alarmante que afecta significativamente la salud y la calidad de vida de las personas \citep{Guthold2020,Swinburn2019Syndemic}. La Organización Mundial de la Salud (OMS) ha identificado el CS como uno de los principales factores de riesgo para enfermedades no transmisibles, tales como enfermedades cardiovasculares, diabetes tipo 2 y ciertos tipos de cáncer \citep{DiPietro2020Advancing,Murray2020GBD,Okunogbe2021Economic,WHO2009GlobalRisks}. Además, el sedentarismo tiene un impacto negativo en la salud mental y genera una disminución en la calidad de vida relacionada con la salud \citep{DiPietro2020Advancing,Guthold2020,Swinburn2019Syndemic}.

Los m├®todos tradicionales de evaluaci├│n del comportamiento sedentario presentan una dicotom├¡a metodol├│gica: los instrumentos subjetivos (cuestionarios de autoinforme) son escalables pero adolecen de sesgos de memoria y deseabilidad social, mientras que las t├®cnicas objetivas de referencia (calorimetr├¡a indirecta, actigraf├¡a de investigaci├│n) se restringen a entornos de laboratorio que no capturan la complejidad del comportamiento en condiciones de vida libre \citep{Healy2024,Prince2008,Migueles2022GRANADA}. Esta limitaci├│n metodol├│gica fundamenta la necesidad de enfoques innovadores que integren objetividad (mediciones continuas mediante sensores), validez ecol├│gica (recopilaci├│n en contextos reales) e interpretabilidad (clasificaciones auditables cl├¡nicamente).

La presente investigaci├│n abord├│ este vac├¡o mediante el desarrollo y validaci├│n de un sistema de inferencia difusa tipo Mamdani que clasifica el comportamiento sedentario a partir de datos biom├®tricos multivariados obtenidos de dispositivos port├ítiles de consumo (Apple Watch) bajo el paradigma \textit{Bring Your Own Device} (BYOD). Este enfoque metodol├│gico integr├│ tres componentes innovadores: (1) normalizaci├│n fisiol├│gica individualizada de variables (Actividad Relativa, Super├ívit Cal├│rico Basal, Delta Card├¡aco, HRV-SDNN), (2) establecimiento de una verdad operativa mediante clustering no supervisado (evitando umbrales heur├¡sticos predefinidos), y (3) validaci├│n rigurosa mediante Leave-One-User-Out para garantizar generalizaci├│n inter-sujeto.

La relevancia cient├¡fica de este trabajo reside en su contribuci├│n al campo emergente de la inteligencia artificial explicable aplicada a salud digital. A diferencia de modelos de caja negra (redes neuronales profundas), el sistema difuso propuesto es inherentemente interpretable mediante reglas ling├╝├¡sticas auditables, una caracter├¡stica cr├¡tica para la adopci├│n cl├¡nica donde la transparencia algor├¡tmica es un imperativo ├®tico. Adem├ís, la demostraci├│n emp├¡rica de que el paradigma BYOD puede generar datos de calidad investigativa sin comprometer la integridad de las mediciones representa un avance significativo hacia la democratizaci├│n de la investigaci├│n en salud, lo cual permite estudios longitudinales a gran escala sin requerir equipamiento especializado costoso.

Desde la perspectiva de salud p├║blica, esta investigaci├│n aporta evidencia sobre la factibilidad de sistemas de vigilancia del sedentarismo basados en datos pasivamente capturados por wearables de consumo masivo. Los resultados obtenidos (F1-Score LOOU = 0.780, validado contra clustering no supervisado) demuestran que es posible clasificar con alta fiabilidad los patrones de comportamiento sedentario en condiciones de vida libre, sentando las bases para el dise├▒o de intervenciones digitales personalizadas y programas de monitoreo comunitario escalables. En un contexto nacional donde la prevalencia de sedentarismo alcanza el 40\% en adultos seg├║n ENSANUT 2022 \citep{Shamah-Levy2023ENSANUT}, la disponibilidad de herramientas objetivas, escalables e interpretables para su evaluaci├│n constituye una contribuci├│n estrat├®gica para el cumplimiento de la meta 3.4 de los Objetivos de Desarrollo Sostenible (reducir en un tercio la mortalidad prematura por ENT para 2030).
